\usepackage[backend=biber, bibstyle=numeric, sorting=none]{biblatex}
\usepackage{tikz}
\usetikzlibrary{math}
\usetikzlibrary{calc}
\usetikzlibrary{fit}
\usetikzlibrary{positioning} %%WARNING: Using the positioning library will render your images unscalable (if no explicitly nodes are present). 
\usetikzlibrary{decorations.pathreplacing}
\usetikzlibrary{decorations.shapes}
\usetikzlibrary{shapes.geometric}
\usetikzlibrary{arrows}
\usepackage{luacode}

%%%%%%%%%%%%%%%%%%%%%%%%%%%%%%%%%%%%%%%%%%%%%%%%%%%%%%%%%%%%%%%
% Start of Luacode
%%%%%%%%%%%%%%%%%%%%%%%%%%%%%%%%%%%%%%%%%%%%%%%%%%%%%%%%%%%%%%%
\begin{luacode*}

	--[[----------------------
	----     Operator     ----
	----------------------]]--

	-- Documentation of the function below
	-- @param originX: the x coordinate of the legend
	-- @param originY: the y coordinate of the legend
	-- @param text: the text of the legend
	-- @param option: the option of the legend
	-- @return: the legend
	function tikzLegend(nargs)
		dico = tikzArgDico(nargs)
		tex.print("\\node["..dico.option.."] at ("..dico.originX..","..dico.originY..") {"..dico.text.."};")
	end

	-- Documentation of the function below
	-- @param option: the option of the legend
	-- @param text: the text of the legend
	-- @param objName: the name of the legend
	-- @return: the legend
	function tikzLegend_relative(nargs)
		dico = tikzArgDico(nargs)
		tex.print("\\node["..dico.option.."] ("..dico.objName..") {"..dico.text.."};")
	end

	-- Documentation of the function below
	-- @param option: the option of the legend
	-- @param text: the text of the legend
	-- @param objName: the name of the legend
	-- @return: the legend
	function tikzLegend_relative2objects(nargs)
		dico = tikzArgDico(nargs)
		local mid_pos = string.format("($(%s)!"..dico.relativeMid.."!(%s)$)", dico.relativeStart, dico.relativeEnd)
		-- local left_pos = string.format("($%s!"..dico.relativeSpace.."!(%s)$)", mid_pos, dico.relativeStart)
		-- local right_pos= string.format("($%s!"..dico.relativeSpace.."!(%s)$)", mid_pos, dico.relativeEnd)
		-- tex.print("\\path["..dico.option.."] "..left_pos.." edge node(testarrow) {"..dico.objName.."}"..right_pos..";")
		tex.print("\\node["..dico.option.."] ("..dico.objName..") at "..mid_pos.." {"..dico.text.."};")
	end

	-- Documentation of the function below
	-- @param originX: the x coordinate of the arrow
	-- @param originY: the y coordinate of the arrow
	-- @param radius: the radius of the arrow
	-- @param option: the option of the arrow
	-- @param objName: the name of the arrow
	-- @return: the arrow
	function tikzArrow_LR(nargs)
	-- default parameters
	  dico = tikzArgDico(nargs)
	  -- draw the arrow
	  local arrow = string.format("\\draw[%s](%s,%s) -- (%s+%s/2,%s) node[above]{%s} -- (%s+%s,%s);",
	  dico.option, dico.originX, dico.originY, dico.originX, dico.radius, dico.originY, dico.objName, dico.originX, dico.radius, dico.originY)
	  tex.print(arrow)
	end

	-- Documentation of the function below
	-- @param originX: the x coordinate of the arrow
	-- @param originY: the y coordinate of the arrow
	-- @param radius: the radius of the arrow
	-- @param option: the option of the arrow
	-- @param objName: the name of the arrow
	-- @return: the arrow
	function tikzArrow_UD(nargs)
	-- default parameters
	dico = tikzArgDico(nargs)
	-- draw the arrow
	local arrow = string.format("\\draw [%s](%s,%s) -- (%s,%s-%s);",
	dico.option, dico.originX, dico.originY, dico.originX, dico.originY, dico.radius, dico.objName, dico.originX, dico.originY, dico.radius)
	local name = string.format("\\node[right] at (%s,%s-%s/2) {%s};",
	dico.originX, dico.originY, dico.radius, dico.objName)
	tex.print(arrow..name)
	end

	-- Documentation of the function below
	-- @param relativeStart: the starting point of the arrow
	-- @param relativeEnd: the ending point of the arrow
	-- @param relativeMid: the middle point of the arrow
	-- @param option: the option of the arrow
	-- @param objName: the name of the arrow
	-- @return: the arrow
	function tikzArrow_relative(nargs)
	-- default parameters
	  dico = tikzArgDico(nargs)
		local mid_pos = string.format("($(%s)!"..dico.relativeMid.."!(%s)$)", dico.relativeStart, dico.relativeEnd)
		local left_pos = string.format("($%s!"..dico.relativeSpace.."!(%s)$)", mid_pos, dico.relativeStart)
		local right_pos= string.format("($%s!"..dico.relativeSpace.."!(%s)$)", mid_pos, dico.relativeEnd)
		tex.print("\\path["..dico.option.."] "..left_pos.." edge node(testarrow) {"..dico.objName.."}"..right_pos..";")
	end
  
	-- Documentation of the function below
	-- @param node: the node on which the arrow is drawn
	-- @param opacity: the opacity of the arrow
	function tikzLineOnCircle(node, opacity)
		--default parameters
		opacity = opacity or 0.5
		-- draw the line
		local string1 = string.format("\\draw[opacity=%s] (%s.north west) -- (%s.south east);", opacity, node, node)
		local string2 = string.format("\\draw[opacity=%s] (%s.north) -- (%s.east);", opacity, node, node)
		local string3 = string.format("\\draw[opacity=%s] (%s.south) -- (%s.west);", opacity, node, node) 
		tex.print(string1..string2..string3)
	end
	-- Documentation of the function below
	-- @param node: the node on which the arrow is drawn
	-- @param opacity: the opacity of the arrow
	function tikzLineOnSquare(nodeS, nodeE, opacity)
		--default parameters
		opacity = opacity or 0.5
		-- draw the line
		local string1 = string.format("\\draw[opacity=%s] (%s.center) -- (%s.center);", opacity, nodeS, nodeE)
		local string2 = string.format("\\draw[opacity=%s] (%s.south) -- (%s.west);", opacity, nodeS, nodeE)
		local string3 = string.format("\\draw[opacity=%s] (%s.east) -- (%s.north);", opacity, nodeS, nodeE) 
		tex.print(string1..string2..string3)
	end
  
	function max(a,b)
	  if a > b then
		return a
	  else
		return b
	  end
	end

	function min(a,b)
	  if a < b then
		return a
	  else
		return b
	  end
	end

	-- Documentation of the function below
	-- @param originX: the x coordinate of the circle
	-- @param originY: the y coordinate of the circle
	-- @param radius: the radius of the circle
	-- return: the minus sign
	function tikzMinus(nargs)
	-- default parameters
	dico = tikzArgDico(nargs)
	-- draw the minus
	local minus = string.format("\\draw [](%s,%s) --(%s+%s,%s);",
	dico.originX, dico.originY, dico.originX, dico.radius, dico.originY)
	tex.print(minus)
	end

	-- Documentation of the function below
	-- @param relativeStart: the starting point of the minus sign
	-- @param relativeEnd: the ending point of the minus sign
	-- return: the minus sign
	function tikzMinus_relative(nargs)
	-- default parameters
	dico = tikzArgDico(nargs)
	-- draw the minus
	local mid_pos = string.format("($(%s)!0.5!(%s)$)", dico.relativeStart, dico.relativeEnd)
	local left_pos = string.format("($%s!%s!(%s)$)", mid_pos, dico.shiftLeft, dico.relativeStart)
	local right_pos= string.format("($%s!%s!(%s)$)", mid_pos, dico.shiftRight, dico.relativeEnd)
	tex.print("\\draw[] "..left_pos.."--"..right_pos..";")
	end

	-- Documentation of the function below
	-- @param originX: the x coordinate of the circle
	-- @param originY: the y coordinate of the circle
	-- @param radius: the radius of the circle
	-- return: the plus sign
	function tikzPlus(nargs)
	-- default parameters
	dico = tikzArgDico(nargs)
	-- draw the plus
	local height = dico.radius/2
	local hori = string.format("\\draw [](%s,%s) --(%s+%s,%s);",
	dico.originX, dico.originY, dico.originX, dico.radius, dico.originY)
	local verti = string.format("\\draw [](%s+%s/2, %s-%s) --(%s+%s/2, %s+%s);",
	dico.originX, dico.radius, dico.originY, height, dico.originX, dico.radius, dico.originY, height)
	tex.print(hori..verti)
	end

	-- Documentation of the function below
	-- @param relativeStart: the starting point of the plus sign
	-- @param relativeEnd: the ending point of the plus sign
	-- return: the plus sign
	function tikzPlus_relative(nargs)
	-- default parameters
	dico = tikzArgDico(nargs)
	-- draw the plus
	local mid_pos = string.format("($(%s)!0.5!(%s)$)", dico.relativeStart, dico.relativeEnd)
	local left_pos = string.format("($%s!%s!(%s)$)", mid_pos, dico.shiftLeft, dico.relativeStart)
	local right_pos= string.format("($%s!%s!(%s)$)", mid_pos, dico.shiftRight, dico.relativeEnd)
	local up_pos = string.format("($%s!%s!90:(%s)$)", mid_pos, dico.shiftUp, dico.relativeEnd)
	local down_pos= string.format("($%s!%s!270:(%s)$)", mid_pos, dico.shiftDown, dico.relativeEnd)
	tex.print("\\draw[] "..mid_pos.."--"..up_pos..";")
	tex.print("\\draw[] "..mid_pos.."--"..down_pos..";")
	tex.print("\\draw[] "..left_pos.."--"..right_pos..";")
	end
	
	-- Documentation of the function below
	-- @param originX: the x coordinate of the circle
	-- @param originY: the y coordinate of the circle
	-- @param radius: the radius of the circle
	-- return: the equal sign
	function tikzEqual(nargs)
	-- default parameters
	dico = tikzArgDico(nargs)
	-- draw the equal
	local height = dico.radius/2
	local upper = string.format("\\draw[](%s,%s+%s) --(%s+%s,%s+%s);",
	dico.originX, dico.originY, height, dico.originX, dico.radius, dico.originY, height)
	local lower = string.format("\\draw[](%s,%s) --(%s+%s,%s);",
	dico.originX, dico.originY, dico.originX, dico.radius, dico.originY)
	tex.print(upper..lower)
	end
	
	-- Documentation of the function below
	-- @param relativeStart: the starting point of the equal sign
	-- @param relativeEnd: the ending point of the equal sign
	-- @param relativeSpace: the space between the equal sign and the start/end point
	-- @param relativeLength: the length of the equal sign
	-- @return: the equal sign
	function tikzEqual_relative(nargs)
	-- default parameters
	dico = tikzArgDico(nargs)
	-- draw the equal
	local mid_pos       = string.format("\\node (mid) at         ($(%s)!%s!(%s)$) {};",   dico.relativeStart, dico.relativeMid, dico.relativeEnd)
	local up_pos        = string.format("\\node (up) at          ($(mid)!%s!90:(%s)$) {};",  dico.relativeSpace, dico.relativeEnd)
	local down_pos      = string.format("\\node (down) at        ($(mid)!%s!270:(%s)$) {};", dico.relativeSpace, dico.relativeEnd)
	local up_left       = string.format("\\node (upleft) at      ($(%s)!%s!90:(mid)$) {};", dico.relativeStart, dico.relativeSpace)
	local up_right      = string.format("\\node (upright) at     ($(%s)!%s!270:(mid)$) {};", dico.relativeEnd, dico.relativeSpace)
	local down_left     = string.format("\\node (downleft) at    ($(%s)!%s!270:(mid)$) {};", dico.relativeStart, dico.relativeSpace)
	local down_right    = string.format("\\node (downright) at   ($(%s)!%s!90:(mid)$) {};", dico.relativeEnd, dico.relativeSpace)
	local uplcorner     = string.format("\\node (uplcorner) at   ($(up)!%s!(upleft)$) {};", dico.relativeLength)
	local uprcorner     = string.format("\\node (uprcorner) at   ($(up)!%s!(upright)$) {};", dico.relativeLength)
	local downlcorner   = string.format("\\node (downlcorner) at ($(down)!%s!(downleft)$) {};", dico.relativeLength)
	local downrcorner   = string.format("\\node (downrcorner) at ($(down)!%s!(downright)$) {};", dico.relativeLength)
	tex.print(mid_pos..up_pos..down_pos..up_left..up_right..down_left..down_right..uplcorner..uprcorner..downlcorner..downrcorner)
	tex.print("\\draw[] (downlcorner)--(downrcorner);")
	tex.print("\\draw[] (uplcorner)--(uprcorner);")
	end

	-- Documentation of the function below
	-- @param relativeStart: the starting point of the equal sign
	-- @param relativeSpace: the space between the equal sign and the start/end point
	-- @param relativeLength: the length of the equal sign
	-- @return: the equal sign
	function tikzEqual_relative1object(nargs)
	-- default parameters
	dico = tikzArgDico(nargs)
	-- draw the equal
	local node_end = string.format("\\node["..dico.option.."] (end) {};", dico.relativeMid, dico.relativeStart)
	local mid_pos       = string.format("\\node (mid) at         ($(%s)!0.5!(end)$) {};",   dico.relativeStart)
	local up_pos        = string.format("\\node (up) at          ($(mid)!%s!90:(end)$) {};",  dico.relativeSpace)
	local down_pos      = string.format("\\node (down) at        ($(mid)!%s!270:(end)$) {};", dico.relativeSpace)
	local up_left       = string.format("\\node (upleft) at      ($(%s)!%s!90:(mid)$) {};", dico.relativeStart, dico.relativeSpace)
	local up_right      = string.format("\\node (upright) at     ($(end)!%s!270:(mid)$) {};", dico.relativeSpace)
	local down_left     = string.format("\\node (downleft) at    ($(%s)!%s!270:(mid)$) {};", dico.relativeStart, dico.relativeSpace)
	local down_right    = string.format("\\node (downright) at   ($(end)!%s!90:(mid)$) {};", dico.relativeSpace)
	local uplcorner     = string.format("\\node (uplcorner) at   ($(up)!%s!(upleft)$) {};", dico.relativeLength)
	local uprcorner     = string.format("\\node (uprcorner) at   ($(up)!%s!(upright)$) {};", dico.relativeLength)
	local downlcorner   = string.format("\\node (downlcorner) at ($(down)!%s!(downleft)$) {};", dico.relativeLength)
	local downrcorner   = string.format("\\node (downrcorner) at ($(down)!%s!(downright)$) {};", dico.relativeLength)
	tex.print(node_end..mid_pos..up_pos..down_pos..up_left..up_right..down_left..down_right..uplcorner..uprcorner..downlcorner..downrcorner)
	tex.print("\\draw[] (downlcorner)--(downrcorner);")
	tex.print("\\draw[] (uplcorner)--(uprcorner);")
	end
  
	function tikzDrawAlongCircle(nb_dots, origin_x, origin_y, radius, angle1, angle2)
	  local function draw_dot_at_angle(angle)
		local x = origin_x + radius * math.cos(math.rad(angle))
		local y = origin_y + radius * math.sin(math.rad(angle))
		tex.sprint(string.format("\\filldraw[black] (%f,%f) circle (%f pt);", x, y, 1))
	  end
	  local angle_step = (angle2 - angle1) / (nb_dots-1)
	  for i=0,nb_dots-1 do
		draw_dot_at_angle(angle1 + i * angle_step)
	  end
	end
  
	function tikzDrawAlongCircle_relative(start_p, end_p)
	  local style = 'ultra thick,dotted, dash pattern=on 2pt off 0.5cm'
	  local string = string.format("\\draw[%s] (%s) to[bend right] (%s);", style, start_p, end_p)
	  tex.print(string)
	end

	-- function tikzDrawAlongCircle_relative(nargs)
	--   dico = tikzArgDico(nargs)
	--   local style = 'ultra thick,dotted, dash pattern=on 2pt off 0.5cm'
	--   local string = string.format("\\draw[%s] (%s) to[bend right] (%s);", style, dico.relativeStart, dico.relativeEnd)
	--   tex.print(string)
	-- end
  
	--[[----------------------
	----    Matrix draw   ----
	----------------------]]--

	-- Documentation of the function below
	-- @param relativeStart: the starting point of the matrix
	-- @param relativeEnd: the ending point of the matrix
	-- @param objName: the name of the matrix
	-- @param objLabel: the label of the matrix
	-- @return: the matrix
	function tikzMatrix2D(nargs)
	  --default parameter
	  dico = tikzArgDico(nargs)
	  --draw the vector
	  local grid = string.format("\\draw ["..dico.optionGrid.."] ("..dico.relativeStart..") node[] ("..dico.objLabelStart..") {} grid ("..dico.relativeEnd..") node[] ("..dico.objLabelEnd..") {"..dico.objName.."};")
	  tex.print(grid)
	end

	-- Documentation of the function below
	-- @param originX: the x coordinate of the matrix
	-- @param originY: the y coordinate of the matrix
	-- @param tensorWidth: the width of the tensor
	-- @param tensorHeight: the height of the tensor
	-- @param tensorColor: the color of the tensor
	-- @param tensorFillColor: the fill color of the tensor
	-- @param objName: the name of the tensor
	-- @param objLabel: the label of the tensor
	-- @return: the tensor
	function tikzFactor(nargs)
	  --default parameter
	  dico = tikzArgDico(nargs)
	  --draw the factor
	  local factor = string.format("\\node[draw, rectangle, minimum width=%fcm, minimum height=%fcm, color=%s, fill=%s] (%s) at (%f,%f) {%s};", dico.tensorWidth, dico.tensorHeight, dico.tensorColor, dico.tensorFillColor, dico.objLabel, dico.originX, dico.originY, dico.objName)
	  tex.print(factor)
	end

	-- Documentation of the function below
	-- @param option: the option of the tensor
	-- @param objLabel: the label of the tensor
	-- @param objName: the name of the tensor
	-- @param tensorWidth: the width of the tensor
	-- @param tensorHeight: the height of the tensor
	-- @param tensorColor: the color of the tensor
	-- @param tensorFillColor: the fill color of the tensor
	-- @return: the tensor
	function tikzFactor_relative(nargs)
	  --default parameter
	  dico = tikzArgDico(nargs)
	  --draw the factor
	  local factor = string.format("\\node[draw, rectangle, minimum width=%fcm, minimum height=%fcm, color=%s, fill=%s, %s] (%s) {%s};", dico.tensorWidth, dico.tensorHeight, dico.tensorColor, dico.tensorFillColor, dico.option, dico.objLabel, dico.objName)
	  tex.print(factor)
	end
  
	--[[----------------------
	----    Tensor draw   ----
	----------------------]]--
  
	-- Documentation of the function below
	-- @param originX: the x coordinate of the tensor
	-- @param originY: the y coordinate of the tensor
	-- @param width: the width of the tensor
	-- @param angle: the angle of the tensor
	-- @param tensorName: the name of the tensor
	-- @param objLabel: the name of the tensor
	-- @return: the tensor
	function tikzCore(nargs)
		--default parameter
		dico = tikzArgDico(nargs)
		--draw the core
		local string = string.format("\\draw[] (%f+%f ,%f) arc (0:%f:%f cm);",
		dico.originX, dico.radius, dico.originY, dico.angle, dico.radius)
		local name = string.format("\\node[] (%s) at (%f,%f) {%s};", dico.objLabel, dico.originX, dico.originY, dico.tensorName)
		tex.print(string..name)
	end
  
	-- Documentation of the function below
	-- @param originX: the x coordinate of the tensor
	-- @param originY: the y coordinate of the tensor
	-- @param angle: the angle of the tensor
	-- @param tensorMaxDim: the maximum dimension of the tensor
	-- @param tensorName: the name of the tensor
	-- @param branchLength: the length of the branch
	-- @param tensorEdgeSize: the size of the edge
	-- @param tensorOutDimName: the name of the output dimension
	-- @return: the tensor with branches
	function tikzCorewBranch(nargs)
		--default parameter
		dico = tikzArgDico(nargs)
		--draw the core
		local angle_seperation = dico.angle/dico.tensorMaxDim;
		tex.print(string.format("\\node[] (core) at (%f,%f) {%s};", dico.originX, dico.originY, dico.tensorName))
		--draw the branches
		for i=1,dico.tensorMaxDim do
			local string = string.format("\\draw[] (%f+%f ,%f ) arc (0:%f:%fcm) node[] (branch"..i..") {};", 
			dico.originX, dico.radius, dico.originY, (i-1)*angle_seperation, dico.radius)
			local branch = string.format("\\draw[] (branch"..i..".center) -- ++(%f:%f) ++(%f:%f) node[] {%s};",
			(i-1)*angle_seperation, dico.branchLength, (i-1)*angle_seperation, dico.tensorEdgeSize, dico.tensorOutDimName[i])
			tex.print(string..branch)
		end
		local string = string.format("\\draw[] (%f+%f ,%f ) arc (0:%f:%fcm) node[] (branch"..dico.tensorMaxDim..") {};",
		dico.originX, dico.radius, dico.originY, (dico.tensorMaxDim)*angle_seperation, dico.radius)
		tex.print(string)
	end

	-- Documentation of the function below
	-- @param originX: the x coordinate of the tensor
	-- @param originY: the y coordinate of the tensor
	-- @param angle: the angle of the tensor
	-- @param tensorMaxDim: the maximum dimension of the tensor
	-- @param tensorName: the name of the tensor
	-- @param tensorTucker: the type of the tensor
	-- @param tensorInnDimName: the name of the input dimension
	-- @param tensorOutDimName: the name of the output dimension
	-- @param tensorEdgeSize: the size of the edge
	-- @param tensorLabel: the label of the tensor
	-- @param startDot: the starting dot
	-- @param endDot: the ending dot
	-- @return: the tensor with branches and dots
	function tikzCoredBranch(nargs)
	 	--default parameter
	 	dico = tikzArgDico(nargs)
		--draw the core
	 	style_factor = ''
	 	if dico.tensorTucker then
	 	  style_factor = 'draw, circle, anchor=center, minimum size=1cm, fill=white'
	 	end
	 	local angle_seperation = dico.angle/dico.tensorMaxDim;
	 	tex.print(string.format("\\node[minimum size = "..dico.radius.."cm, label = center : "..dico.tensorName.."] ("..dico.tensorLabel..") at (%f,%f) {};", dico.originX, dico.originY))
	 	--draw the branches
	 	for i=1,dico.tensorMaxDim do
	 		local string = string.format("\\draw[] ("..dico.originX.."+"..dico.radius..","..dico.originY..") arc (0:"..(i-1)*angle_seperation..":"..dico.radius.."cm) node[] (mid"..i..") {};")
			tex.print(string)
			if dico.tensorTucker then
				local branch = string.format("\\draw[] (mid"..i..".center) -- ++("..(i-1)*angle_seperation..":"..dico.tensorEdgeSize..") node[minimum size = "..dico.tensorEdgeSize.."cm, label = center : "..dico.tensorInnDimName[i]..","..style_factor.."] (dBranch"..i..") {};")
				tex.print(branch)
			else
				local branch = string.format("\\draw[] (mid"..i..".center) -- ++("..(i-1)*angle_seperation..":"..dico.tensorEdgeSize..") ++("..(i-1)*angle_seperation..":"..dico.tensorEdgeSize..") node[minimum size = "..dico.tensorEdgeSize.."cm, label = center : "..dico.tensorInnDimName[i]..","..style_factor.."] (dBranch"..i..") {};")
				tex.print(branch)
			end
	 		local string_dbranch = string.format("\\node (res"..i..") at ($(mid"..i..")!0.5!(dBranch"..i..")$) {};")
			tex.print(string_dbranch)
			if dico.tensorTucker then
			  local branch = string.format("\\draw[] (dBranch"..i..") -- ++("..(i-1)*angle_seperation..":"..dico.tensorEdgeSize.."/2) ++("..(i-1)*angle_seperation..":"..dico.tensorEdgeSize.."/2) node[] (end"..i..") {"..dico.tensorOutDimName[i].."};")
			  tex.print(branch)
			end
	 	end
		local string = string.format("\\draw[] (%f+%f ,%f ) arc (0:%f:%fcm) node[] ("..dico.tensorInnDimName[dico.tensorMaxDim]..") {};",
		dico.originX, dico.radius, dico.originY, (dico.tensorMaxDim)*angle_seperation, dico.radius)
		tex.print(string)
		tikzDrawAlongCircle_relative(string.format("res%s",dico.startDot), string.format("res%s",dico.endDot))
	end
  
	-- Documentation of the function below
	-- @param originX: the x coordinate of the tensor
	-- @param originY: the y coordinate of the tensor
	-- @param scale: the scale of the tensor
	-- @param tensorMaxDim: the maximum dimension of the tensor
	-- @param tensorCoreMinSize: the minimum size of the core
	-- @param tensorLevelDistance: the distance between levels
	-- @param tensorSiblingDistance: the distance between siblings
	-- @param tensorInnDimName: the name of the input dimension
	-- @param tensorOutDimName: the name of the output dimension
	-- @return: the Tucker tensor
	function tikzTuckerTTN(nargs)
	  --default parameter
	  dico = tikzArgDico(nargs)
	  --draw the Tucker
	  tex.print("\\tikzset{scale = "..dico.scale..", semithick, minimum size="..dico.tensorCoreMinSize.."cm, level distance = "..dico.tensorLevelDistance.."cm, sibling distance="..dico.tensorSiblingDistance.."cm, split/.style = {level distance=}}")
	  --Tree
	  local root = string.format("\\node[circle,draw,minimum size="..1.5*dico.tensorCoreMinSize.."cm] (root) at ("..dico.originX..","..dico.originY..") {"..dico.tensorInnDimName[1].."}")
	  for i = 1,dico.tensorMaxDim,1 do 
		root = root..string.format("child { node[circle,draw] (u"..i..") {"..dico.tensorInnDimName[i+1].."} child { node[] (n"..i..") {"..dico.tensorOutDimName[i].."} } edge from parent node[sloped,anchor=south,auto=false,above,draw=none]{"..dico.tensorInnDimName[i].."} }")
	  end
	  root = root..";"
	  tex.print(root)
	  
	  tex.print("\\tikzset{scale = 1};")
	end

	-- Documentation of the function below
	-- @param originX: the x coordinate of the tensor
	-- @param originY: the y coordinate of the tensor
	-- @param scale: the scale of the tensor
	-- @param tensorMaxDim: the maximum dimension of the tensor
	-- @param tensorCoreMinSize: the minimum size of the core
	-- @param tensorLevelDistance: the distance between levels
	-- @param tensorSiblingDistance: the distance between siblings
	-- @param tensorInnDimName: the name of the input dimension
	-- @param tensorOutDimName: the name of the output dimension
	-- @param treeSiblingDistance: the distance between siblings
	-- @param tensorName: the name of the tensor
	-- @return: the Tucker tensor
	function tikzTensorTrain(nargs)
		--default parameter
		dico = tikzArgDico(nargs)
		-- Set style
		tex.print("\\tikzset{scale = "..dico.scale..", sibling distance = "..dico.treeSiblingDistance.."cm, level distance = "..dico.treeLevelDistance.."cm, split/.style = {level distance=}}")
		-- Tree
		local tree = string.format("\\node[circle, draw, minimum size = "..dico.tensorCoreMinSize.."cm, label = center : "..dico.tensorName[1].."] (root) at ("..dico.originX..","..dico.originY..") {} child[grow="..dico.tensorInnDimOrientation.."] {node (n1) {"..dico.tensorOutDimName[1].."}} child[grow="..dico.tensorOutDimOrientation.."]{");
		for i = 2,dico.tensorMaxDim-1,1 do
		   	if dico.ttDots and i == (dico.tensorMaxDim-1) then
				local mid_dot = dico.treeLevelDistance/2
			   tree = tree..string.format("node[minimum size = "..dico.tensorCoreMinSize.."cm, label = center : "..dico.tensorName[i].."] (node"..i..") {} child[grow="..dico.tensorInnDimOrientation..", level distance="..mid_dot.."cm, white] {node[black] (n"..i..") {"..dico.tensorOutDimName[i].."}} child[grow="..dico.tensorOutDimOrientation.."]{")
		   	else
		   		tree = tree..string.format("node[circle, draw, minimum size = "..dico.tensorCoreMinSize.."cm, label = center : "..dico.tensorName[i].."] (node"..i..") {} child[grow="..dico.tensorInnDimOrientation.."] {node[] (n"..i..") {"..dico.tensorOutDimName[i].."}} child[grow="..dico.tensorOutDimOrientation.."]{")
			end
		end
		tree = tree..string.format("node[circle, draw, minimum size = "..dico.tensorCoreMinSize.."cm, label = center : "..dico.tensorName[dico.tensorMaxDim].."] (node"..dico.tensorMaxDim..") {} child[grow="..dico.tensorInnDimOrientation.."] {node[] (n"..dico.tensorMaxDim..") {"..dico.tensorOutDimName[dico.tensorMaxDim].."} }")
		for i = 2,dico.tensorMaxDim-1,1 do
		  tree = tree..string.format("edge from parent node[sloped,anchor=south,auto=false,above,draw=none]{"..dico.tensorInnDimName[dico.tensorMaxDim+1-i].."}}")
		end
		tree = tree.."edge from parent node[sloped,anchor=south,auto=false,above,draw=none]{"..dico.tensorInnDimName[1].."}};"
		tex.print(tree)
		tex.print("\\tikzset{scale = 1};");
	end
	
	function tikzTensorTrainNoChild(nargs)
		--default parameter
		dico = tikzArgDico(nargs)
		-- Set style
		tex.print("\\tikzset{scale = "..dico.scale..", sibling distance = "..dico.treeSiblingDistance.."cm, level distance = "..dico.treeLevelDistance.."cm, split/.style = {level distance=}}")
		-- Tree
		local tree = string.format("\\node[circle, draw, minimum size = "..dico.tensorCoreMinSize.."cm, label = center : "..dico.tensorName[1].."] (root) at ("..dico.originX..","..dico.originY..") {} child[missing] child[grow="..dico.tensorOutDimOrientation.."]{");
		for i = 2,dico.tensorMaxDim-1,1 do
		   	if dico.ttDots and i == (dico.tensorMaxDim-1) then
				local mid_dot = dico.treeLevelDistance/2
			   tree = tree..string.format("node[minimum size = "..dico.tensorCoreMinSize.."cm, label = center : "..dico.tensorName[i].."] (node"..i..") {} child[missing] child[grow="..dico.tensorOutDimOrientation.."]{")
		   	else
		   		tree = tree..string.format("node[circle, draw, minimum size = "..dico.tensorCoreMinSize.."cm, label = center : "..dico.tensorName[i].."] (node"..i..") {} child[missing] child[grow="..dico.tensorOutDimOrientation.."]{")
			end
		end
		tree = tree..string.format("node[circle, draw, minimum size = "..dico.tensorCoreMinSize.."cm, label = center : "..dico.tensorName[dico.tensorMaxDim].."] (node"..dico.tensorMaxDim..") {} child[missing]")
		for i = 2,dico.tensorMaxDim-1,1 do
		  tree = tree..string.format("edge from parent node[sloped,anchor=south,auto=false,above,draw=none]{"..dico.tensorInnDimName[dico.tensorMaxDim+1-i].."}}")
		end
		tree = tree.."edge from parent node[sloped,anchor=south,auto=false,above,draw=none]{"..dico.tensorInnDimName[1].."}};"
		tex.print(tree)
		tex.print("\\tikzset{scale = 1};");
	end

	-- Documentation of the function below
	-- @param originX: the x coordinate of the tensor
	-- @param originY: the y coordinate of the tensor
	-- @param angle: the angle of the tensor
	-- @param tensorMaxDim: the maximum dimension of the tensor
	-- @param tensorName: the name of the tensor
	-- @param tensorTucker: the type of the tensor
	-- @param tensorInnDimName: the name of the input dimension
	-- @param tensorOutDimName: the name of the output dimension
	-- @param tensorEdgeSize: the size of the edge
	-- @param tensorLabel: the label of the tensor
	-- @param startDot: the starting dot
	-- @param endDot: the ending dot
	-- @return: the tensor with branches and dots
	function tikzTucker(nargs)
	  --default parameter
	  --dico = tikzArgDico(nargs)
	  --draw the core
	  tikzCoredBranch(setmetatable(nargs, {__index = {tensorTucker = true}}))
	end

	function tikzDrawCube(nargs)
		--default parameter
		dico = tikzArgDico(nargs)
		--draw the cube
		local cube = string.format("\\pic ["..dico.color.."] at ("..dico.originX..","..dico.originY..","..dico.originZ..") {"..dico.option.."};")
		local node = string.format("\\node[] ("..dico.objLabel..") at ("..dico.originX..","..dico.originY..","..dico.originZ..") {"..dico.text.."};")
		tex.print(cube..node)
	end

	function tikzTensorTrainNoChildCustom(nargs)
		--default parameter
		dico = tikzArgDico(nargs)
		-- Set style
		tex.print("\\tikzset{scale = "..dico.scale..", sibling distance = "..dico.treeSiblingDistance.."cm, level distance = "..dico.treeLevelDistance.."cm, split/.style = {level distance=}}")
		-- Tree
		local tree = string.format("\\node[ minimum size = "..dico.tensorCoreMinSize.."cm, label = center : "..dico.tensorName[1].."] (root) at ("..dico.originX..","..dico.originY..") {} child[missing] child[grow="..dico.tensorOutDimOrientation.."]{");
		for i = 2,dico.tensorMaxDim-1,1 do
		   	if dico.ttDots and i == (dico.tensorMaxDim-1) then
				local mid_dot = dico.treeLevelDistance/2
			   tree = tree..string.format("node[minimum size = "..dico.tensorCoreMinSize.."cm, label = center : "..dico.tensorName[i].."] (node"..i..") {} child[missing] child[grow="..dico.tensorOutDimOrientation.."]{")
		   	else
		   		tree = tree..string.format("node[  minimum size = "..dico.tensorCoreMinSize.."cm, label = center : "..dico.tensorName[i].."] (node"..i..") {} child[missing] child[grow="..dico.tensorOutDimOrientation.."]{")
			end
		end
		tree = tree..string.format("node[  minimum size = "..dico.tensorCoreMinSize.."cm, label = center : "..dico.tensorName[dico.tensorMaxDim].."] (node"..dico.tensorMaxDim..") {} child[missing]")
		for i = 2,dico.tensorMaxDim-1,1 do
		  tree = tree..string.format("edge from parent node[sloped,anchor=south,auto=false,above,draw=none]{"..dico.tensorInnDimName[dico.tensorMaxDim+1-i].."} edge from parent[draw=none]}")
		end
		tree = tree.."edge from parent node[sloped,anchor=south,auto=false,above,draw=none]{"..dico.tensorInnDimName[1].."} edge from parent[draw=none]};"
		tex.print(tree)
		tex.print("\\tikzset{scale = 1};");
	end

  function tikzTridiagColored(nargs)
  --default parameter
  dico = tikzArgDico(nargs)
  -- Set style
  local grid = string.format("");
  local colorL = 3;
  local startColor = dico.originX;
  local endColor = 0;
  for j = 1,dico.maxBranch,1 do
    startColor = max(dico.endX - j + 2 - colorL,dico.originX);
    if j == dico.maxBranch then
      colorL = 2;
    end
    endColor = min(startColor + colorL, dico.endX);
    --Before color
    grid = grid..string.format("\\draw ["..dico.optionGrid.."] ("..dico.originX..","..j..") node[] ("..dico.objLabelStart..") {} grid ("..startColor..","..j.."+1) node[] ("..dico.objLabelEnd..") {};")
    --After color
    grid = grid..string.format("\\draw ["..dico.optionGrid.."] ("..endColor..","..j..") node[] ("..dico.objLabelStart..") {} grid ("..dico.endX..","..j.."+1) node[] ("..dico.objLabelEnd..") {"..dico.objName.."};")
    -- Colored Block
    grid = grid..string.format("\\draw ["..dico.color.."] ("..startColor..","..j..") node[] ("..dico.objLabelStart..") {} grid ("..endColor..","..j.."+1) node[] ("..dico.objLabelEnd..") {"..dico.objName.."};")
  end
	tex.print(grid)
  end

  function tikzOneBlock(nargs)
	--default parameter
	dico = tikzArgDico(nargs)
	-- Set style
	if dico.linecolor == {} then
		dico.linecolor = 'black';
  	end
  	tex.print("\\draw [step="..dico.step..",draw="..dico.linecolor..",fill="..dico.color.." ] "..dico.relativeStart.." node[] (blockStart) {} grid "..dico.relativeStart.."+"..dico.relativeMid.." node[] (blockEnd) {} rectangle node[] {} "..dico.relativeStart..";")
  end

  function tikzOneRowColored(nargs)
  --default parameter
  dico = tikzArgDico(nargs)
  -- Set style
  -- Set line color if not enough in space
  if #dico.linecolor < dico.endX-dico.originX then
  	for i = #dico.linecolor,dico.endX-dico.originX+1,1 do
		dico.linecolor[i] = 'black';
	end
  end
  if #dico.caseName < dico.endX-dico.originX then
  	for i = #dico.caseName,dico.endX-dico.originX+1,1 do
		dico.caseName[i] = '';
	end
  end
  --main code
  local grid = string.format("");
  local j = 1;
  for i = dico.originX,dico.endX-1,1 do
    grid = grid..string.format("\\draw [draw="..dico.linecolor[j]..",fill="..dico.color[j].." ] ("..i..","..dico.originY..") node[] ("..dico.objLabelStart..") {} grid ("..i.."+1,"..dico.originY.."+1) node[] ("..dico.objLabelEnd..") {"..dico.objName.."} rectangle node[] {"..dico.caseName[j].."} ("..i..","..dico.originY..");")
    j = j + 1;
  end
	tex.print(grid)
  end

  function tikzOneColumnColored(nargs)
  --default parameter
  dico = tikzArgDico(nargs)
  -- Set style
  -- Set line color if not enough in space
  if #dico.linecolor < dico.endY-dico.originY then
  	for i = #dico.linecolor,dico.endY-dico.originY+1,1 do
		dico.linecolor[i] = 'black';
	end
  end
  if #dico.caseName < dico.endY-dico.originY then
  	for i = #dico.caseName,dico.endY-dico.originY+1,1 do
		dico.caseName[i] = '';
	end
  end
  --main code
  local grid = string.format("");
  local j = 1;
  for i = dico.originY,dico.endY-1,1 do
    grid = grid..string.format("\\draw [step="..dico.step..",draw="..dico.linecolor[j]..",fill="..dico.color[j]..", "..dico.option.."] ("..dico.originX..","..i..") node[] ("..dico.objLabelStart..") {} grid ("..dico.originX.."+1,"..i.."+1) node[] ("..dico.objLabelEnd..") {"..dico.objName.."} rectangle node[] {"..dico.caseName[j].."} ("..dico.originX..","..i..");")
    j = j + 1;
  end
	tex.print(grid)
  end

	function tikzArgDico(dico)
		setmetatable(dico, {__index = {originX=0}})
		-- Standard dico
		 dico.originX 		 	= dico.originX or 0
	  	 dico.originY 		 	= dico.originY or 0
		 dico.originZ 		 	= dico.originZ or 0
		 dico.endX 			 	= dico.endX or 0
		 dico.endY 			 	= dico.endY or 0
		 dico.endZ 			 	= dico.endZ or 0
	  	 dico.radius 		 	= dico.radius or 1
		 dico.color 		 	= dico.color or 'black'
	  	 dico.angle 		 	= dico.angle or 360
	  	 dico.maxBranch 	 	= dico.maxBranch or 1
	  	 dico.outerEdgesLength  = dico.outerEdgesLength or 0
		 dico.text 			 	= dico.text or ''
		 dico.option 		 	= dico.option or ''
		 dico.objName 		 	= dico.objName or ''
		 dico.objLabel 	 		= dico.objLabel or ''
		 dico.objLabelStart 	= dico.objLabelStart or ''
		 dico.objLabelEnd 	 	= dico.objLabelEnd or ''
		 dico.scale 		 	= dico.scale or 1
		 dico.optionGrid 		= dico.optionGrid or 'step=1.0, blue, very thick'
		 dico.linecolor 		= dico.linecolor or {}
		 dico.caseName 		 	= dico.caseName or {}
		 dico.shiftLeft 	 	= dico.shiftLeft or 0
		 dico.shiftRight 	 	= dico.shiftRight or 0
		 dico.shiftUp 		 	= dico.shiftUp or 0.5
		 dico.shiftDown 	 	= dico.shiftDown or 0.5
		 dico.step 			 	= dico.step or 1
		-- Relative dico
		 dico.relativeStart 	= dico.relativeStart or ''
		 dico.relativeEnd 	 	= dico.relativeEnd or ''
		 dico.relativeMid 	 	= dico.relativeMid or 0.3
		 dico.relativeSpace 	= dico.relativeSpace or 0.5
		 dico.relativeLength 	= dico.relativeLength or 0.3
		-- Tensor dico
		 dico.tensorName 	 	= dico.tensorName or ''
		 dico.tensorLabel 	 	= dico.tensorLabel or ''
		 dico.tensorColor 	 	= dico.tensorColor or 'black'
		 dico.tensorFillColor 	= dico.tensorFillColor or 'white'
		 dico.tensorWidth 	 	= dico.tensorWidth or 1
		 dico.tensorHeight 	 	= dico.tensorHeight or 1
		 dico.tensorCoreMinSize = dico.tensorCoreMinSize or 1
		 dico.tensorMaxDim 	 	= dico.tensorMaxDim or 1
		 dico.tensorOutDimName 	= dico.tensorOutDimName or ''
		 dico.tensorOutDimLabel = dico.tensorOutDimLabel or ''
		 dico.tensorInnDimName 	= dico.tensorInnDimName or ''
		 dico.tensorInnDimLabel = dico.tensorInnDimLabel or ''
		 dico.tensorEdgeSize 	= dico.tensorEdgeSize or 1
		 dico.tensorTucker 	 	= dico.tensorTucker or false
		-- Tensor Train dico
		 dico.tensorInnDimOrientation = dico.tensorInnDimOrientation or 'down'
		 dico.tensorOutDimOrientation = dico.tensorOutDimOrientation or 'right'
		 dico.ttDots 				  = dico.ttDots or false
		-- Tree dico
		 dico.treeLevelDistance   = dico.treeLevelDistance or 1.5
		 dico.treeSiblingDistance = dico.treeSiblingDistance or 2
		-- Dot between two nodes
		 dico.startDot = dico.startDot or 1
		 dico.endDot   = dico.endDot or 1

	  	return dico
	end

	function tikzShowDico(dico)
	local i =0;
    for k,v in pairs(dico) do
      tex.print("\\node[] at (0,"..i..") {"..k.."};")
	  if type(v) == "table" then
		for k1,v1 in pairs(v) do
			tex.print("\\node[] at (5,"..i..") {"..k1.."};")
			tex.print("\\node[] at (10,"..i..") {"..tostring(v1).."};")
			i = i + 1;
		end
	  else
      		tex.print("\\node[] at (10,"..i..") {"..tostring(v).."};")
      		i = i + 1;
	  end
    end
	end
  \end{luacode*}
  